\SourceLeaf{254}%
verwenden wir auch hier. Für einen beliebigen Punkt der Fläche sei \(r\) der kürzeste Abstand von einem festen Punkt, und \(\varphi\) sei der Winkel an diesem festen Punkt zwischen dem ersten Element von \(r\) und einer festen Richtung. Damit gilt \(E'=1\), \(F'=0\), \(\omega'=90^0\). Außerdem setzen wir \(\BareRoot{G'}=m\), sodass ein beliebiges Linienelement \(=\BareRoot{(\mathrm{d}r^2+mm\mathrm{d}\varphi^2)}\) ist. Die vier im vorhergehenden Artikel für \(\alpha,\beta,\gamma,\delta\) hergeleiteten Gleichungen ergeben somit:
\[
\begin{aligned}
\BareRoot{E}\mathord{.}\cos(\omega-\psi)&=\frac{\mathrm{d}r}{\mathrm{d}p}\quad .\ .\ .\ .\ .\ .\ .\ .\ .\ .\ .\ .\quad(1)\\
\BareRoot{G}\mathord{.}\cos\psi&=\frac{\mathrm{d}r}{\mathrm{d}q}\quad .\ .\ .\ .\ .\ .\ .\ .\ .\ .\ .\ .\quad(2)\\
\BareRoot{E}\mathord{.}\sin(\psi-\omega)&=m\mathord{.}\frac{\mathrm{d}\varphi}{\mathrm{d}p}\quad .\ .\ .\ .\ .\ .\ .\ .\ .\ .\ .\quad(3)\\
\BareRoot{G}\mathord{.}\sin\psi&=m\mathord{.}\frac{\mathrm{d}\varphi}{\mathrm{d}q}\quad .\ .\ .\ .\ .\ .\ .\ .\ .\ .\ .\quad(4)
\end{aligned}
\]
Die letzte und die vorletzte Gleichung jenes Artikels ergeben dagegen
\[
\begin{aligned}
EG-FF&=E(\frac{\mathrm{d}r}{\mathrm{d}q})^2-2F\mathord{.}\frac{\mathrm{d}r}{\mathrm{d}p}\mathord{.}\frac{\mathrm{d}r}{\mathrm{d}q}+G(\frac{\mathrm{d}r}{\mathrm{d}p})^2\quad .\ .\ .\ .\ .\ .\quad(5)\\
(E\mathord{.}\frac{\mathrm{d}r}{\mathrm{d}q}-F\mathord{.}\frac{\mathrm{d}r}{\mathrm{d}p})\mathord{.}\frac{\mathrm{d}\varphi}{\mathrm{d}q}&=(F\mathord{.}\frac{\mathrm{d}r}{\mathrm{d}q}-G\mathord{.}\frac{\mathrm{d}r}{\mathrm{d}p})\mathord{.}\frac{\mathrm{d}\varphi}{\mathrm{d}p}\quad .\ .\ .\ .\ .\ .\ .\quad(6)
\end{aligned}
\]

Aus diesen Gleichungen sind die Größen \(r,\varphi,\psi\) und, falls es nötig erscheint, \(m\) aus \(p\) und \(q\) zu bestimmen. Die Integration der Gleichung (5) liefert \(r\). Ist diese Größe gefunden, so liefert die Integration der Gleichung (6) \(\varphi\), und eine der beiden Gleichungen (1) und (2) liefert \(\psi\) selbst. Schließlich ergibt sich \(m\) aus einer der beiden Gleichungen (3) und (4).

Bei der allgemeinen Integration der Gleichungen (5) und (6) müssen zwei willkürliche Funktionen auftreten. Ihre Bedeutung ist leicht zu verstehen, wenn man bedenkt, dass diese Gleichungen über den hier betrachteten Fall hinaus ebenso gelten, wenn \(r\) und \(\varphi\) im allgemeineren Sinn von Artikel 16 aufgefasst werden. Dann ist \(r\) die Länge einer kürzesten Linie, die senkrecht zu einer fest vorgegebenen, beliebig gewählten Linie gezogen wird, und \(\varphi\) eine willkürliche Funktion der Länge des Abschnitts dieser Linie zwischen einer veränderlichen kürzesten Linie und einem beliebig gewählten festen Punkt. Die allgemeine Lösung muss daher all diese Möglichkeiten umfassen. Die willkürlichen Funktionen werden erst dann bestimmt, wenn jene beliebig gewählte Linie und die Funktion ihrer Abschnitte, die \(\varphi\) darstellen soll, vorgegeben sind. In unserem Fall können wir einen unendlich kleinen Kreis wählen, dessen Mittelpunkt der Punkt ist, von dem aus die Abstände \(r\) gemessen werden. Dann bezeichnet \(\varphi\) die Längen der Bögen dieses Kreises, geteilt durch den Radius,
